\subsection{An ordinary source cannot be monomial}
\label{sec:kt-ordinary-imprimitive-exclusion}

This is a further source-scope reduction.  It gives no bound on
the strict Selmer group and is not needed for the determinant
normalization or self-dual exclusion above.

\begin{lemma}[KT-ORDINARY-IMPRIMITIVE-EXCLUSION]
Let $p>3$, let $E$ be a number field unramified at a place
$v\mid p$, and let $e>0$ with $(p-1)\nmid e$.
Let $R:G_E\to\mathrm{GL}_3(K)$ be a continuous absolutely
irreducible representation over a finite extension $K$ of
$\mathbf Q_p$.  Suppose
\[
              \bar R^{\mathrm{ss}}=1\oplus\omega^e
                                            \oplus\omega^{-e}
                                                               \tag{IM.1}
\]
and $R|_{G_{E_v}}$ has a stable full flag whose inertia
characters are $(\varepsilon_p^e,1,\varepsilon_p^{-e})$.
Then $R$ does not preserve a set of three lines with direct sum
equal to its underlying space, even after a finite coefficient
extension.  Its projective image is infinite.
\end{lemma}

\begin{proof}
Suppose first that three such lines $L_1,L_2,L_3$ exist.
Their permutation action is transitive, since a union of any
proper collection of orbits would span a proper nonzero
$G_E$-subrepresentation.  Let $P\subset S_3$ be the permutation
image and $H$ its kernel.

\paragraph{Reduction retains the permutation quotient.}
There is a stable monomial lattice.  Choose $0\ne v_1\in L_1$
and choose $g_i$ carrying $L_1$ to $L_i$; set $v_i=R(g_i)v_1$.
The stabilizer of $L_1$ acts on it by a continuous character
of a compact group, hence by units of the coefficient field.
For any $g\in G_E$, write $g g_i=g_j h$ with $h$ in that
stabilizer.  Then $R(g)v_i$ is a unit multiple of $v_j$.
Consequently $\bigoplus_i\mathcal O_K v_i$ is stable, after
one finite coefficient extension defining the lines.
Its reduction consists of monomial matrices with the same
permutation quotient $P$: the nonzero entries remain units.

By (IM.1), the finite residual image can also be triangularized.
The kernel of its diagonal-character map is a normal
$p$-subgroup, and the quotient is abelian.  Its image in $P$
is trivial because $|P|$ divides six and $p>3$.
Thus $P$ is abelian.  The only abelian transitive subgroup
of $S_3$ is $C_3$, so
\[
                          G_E/H\simeq C_3.                \tag{IM.2}
\]

\paragraph{The permutation action is trivial on $p$-adic inertia.}
Put $I=I_v$ and $I_0=I\cap H$.  This is a finite-index normal
subgroup of $I$.  The three lines are $I_0$-stable, and their
three characters are, in some order,
\[
                  \varepsilon_p^e|_{I_0},\quad1,
                              \varepsilon_p^{-e}|_{I_0}.  \tag{IM.3}
\]
Indeed these are the composition characters of the restricted
ordinary flag.  They are pairwise distinct: the cyclotomic
image on $I_0$ remains infinite, whereas $e>0$.
Each character in (IM.3) is invariant under conjugation by $I$.
Conjugation must therefore preserve each of the three distinct
character lines, so $I\subset H$.

\paragraph{The cubic self-twist contradicts the residual characters.}
Choose a nontrivial character
$\tau:G_E/H\to\mu_3$.  The monomial representation is induced
from the stabilizer $H$ of one line, and hence
\[
                              R\otimes\tau\simeq R.        \tag{IM.4}
\]
One can see this without an induction theorem: on the above
basis put $T v_i=\tau(g_i)v_i$.  The identity
$\tau(g)\tau(g_i)=\tau(g_j)$ gives
$T R(g)T^{-1}=\tau(g)R(g)$.
Since $I\subset H$, $\tau$ is unramified at $v$.
Reduction is injective on $\mu_3$ because $p>3$, so
$\bar\tau$ is still nontrivial.

Reducing (IM.4) and using (IM.1) gives equality of global
character multisets
\[
 \{1,\omega^e,\omega^{-e}\}
       =\{\bar\tau,\bar\tau\omega^e,
                                  \bar\tau\omega^{-e}\}. \tag{IM.5}
\]
In particular $\bar\tau=\omega^e$ or $\omega^{-e}$.
Both are nontrivial on $I_v$: $E_v/\mathbf Q_p$ is unramified
and $(p-1)\nmid e$.  This contradicts the unramifiedness of
$\tau$ at $v$.  This also covers the case in which
$\omega^e$ has order three.

Finally, if the projective image of $R$ were finite, a
finite-index subgroup $I'\subset I_v$ would act by scalar
matrices.  Its three composition characters would be equal.
The ordinary flag gives instead
$\varepsilon_p^e|_{I'},1,\varepsilon_p^{-e}|_{I'}$, which
are pairwise distinct as above.  This is impossible.
\end{proof}

For the actual source, the hypotheses on the field and exponent
follow from GD, and the ordered flag follows from
\texttt{LevelAt} at modulus two.  The lemma uses no auxiliary
unit condition and does not assume that a characteristic-zero
ordinary source can be recovered from an ambient residual
Hecke eigenvalue.  Combined with the three-unit exclusion,
it leaves primitive, nonessentially-self-dual actual stable
systems with infinite projective image.  No theorem comparing
their integral deformation rings with a Hecke algebra is
deduced here.
